% Contact Nuno Laranjeiro cnl@dei.uc.pt for any feedback, suggestions, or corrections regarding this template.

% ************************************ CONFIGURATION *****************************************
% Fill in _all_ items in this configuration section with your Thesis details.                                                                      

% If you are writing a Phd Thesis, change *msc* to *phd* in the following line
\def\thesistype{msc}                                                     

% docLanguage is the main language of the document. Possible options are: en,pt
\def\docLanguage{en}

% Tentative title -- revisit after the literature review (Section 3.5) fixes the research questions.
\def\titleEN{Towards Fair Algorithmic Recourse with Kolmogorov-Arnold Networks}
\def\subtitleEN{}

\def\titlePT{Rumo a um Recurso Algorítmico Justo com Redes de Kolmogorov-Arnold}
\def\subtitlePT{}

% MSc Theses - cover page
\def\descMScEN{Dissertation in the context of the Master in Informatics Engineering, specialization in Artificial Intelligence, advised by
Professor Pedro Henriques Abreu and Professor Bruno Veloso and presented to the Department of Informatics Engineering of the Faculty of Sciences and Technology of the University of Coimbra.}

\def\descMScPT{Dissertação no âmbito do Mestrado em Engenharia Informática, especialização em Inteligência Artificial, orientada
pelo Professor Pedro Henriques Abreu e pelo Professor Bruno Veloso e apresentada ao Departamento de Engenharia Informática da Faculdade de Ciências e Tecnologia da Universidade de Coimbra.}

% PhD Theses - cover page
\def\descPhDEN{Thesis in the context of the <full name of the doctoral program, including branches/areas if applicable>, advised by Professor 
<full name of the supervisor> and presented at <name of the academic unit>/<name of the department within the academic unit, if applicable>.}

\def\descPhDPT{Tese no âmbito do <nome completo do doutoramento, incluindo ramos/área se aplicável> orientada pelo/a Professor/a <nome completo 
do orientador> e apresentada <à nome da unidade orgânica>/<ao nome do departamento, se aplicável da/do nome da unidade orgânica>.}


% MSc Theses - title Page
\def\titlePageDescMScEN{Dissertation in the context of the Master in Informatics Engineering, specialization in Artificial Intelligence, advised by Professor
Pedro Henriques Abreu and Professor Bruno Veloso and presented to the Department of Informatics Engineering of the Faculty of Sciences and Technology of the University of Coimbra.}
\def\titlePageDescMScPT{Dissertação no âmbito do Mestrado em Engenharia Informática, especialização em Inteligência Artificial, orientada
pelo Professor Pedro Henriques Abreu e pelo Professor Bruno Veloso e apresentada ao Departamento de Engenharia Informática da Faculdade de Ciências e Tecnologia da Universidade de Coimbra.}


% PhD Theses - title Page
\def\titlePageDescPhdEN{PhD Thesis submitted to the University of Coimbra\\Advised by Professor <Supervisor's name>.}
\def\titlePageDescPhdPT{Tese de Doutoramento submetida à Universidade de Coimbra\\Orientada pelo Professor <Nome do Orientador>.}


\def\author{Tiago Silva}
\def\dateEN{January 2027}
\def\datePT{Janeiro 2027}
% ********************************** END CONFIGURATION ***************************************



\documentclass[twoside, 12pt, a4paper]{book}
\raggedbottom

\usepackage[a4paper,top=2.5cm,bottom=2.5cm,left=3.5cm,right=2.5cm,headsep=12pt]{geometry}
\usepackage[final]{pdfpages}

\usepackage[absolute,overlay]{textpos}


\usepackage{ifthen}

\usepackage[T1]{fontenc}
\usepackage[utf8]{inputenc}

\ifthenelse{\equal{\docLanguage}{en}}
{
	\usepackage[english]{babel}
}
{
	\usepackage[portuguese]{babel}
}

\usepackage{graphicx}
\usepackage{glossaries}
\usepackage{tabularx}
\usepackage{mathpazo}
\usepackage{xcolor, soulutf8, afterpage}  
\usepackage{emptypage}
\usepackage{setspace}

\usepackage[disable]{todonotes}

\ifthenelse{\equal{\thesistype}{msc}}
{

	%Cor do fundo da capa, lombada e contracapa: C 17%, M 27%, Y 45%, K 4% (~ Pantone 7753C)
	\definecolor{cover-background}{cmyk}{0.17 0.27 0.45 0.04}
	%Cor das marcas e Tipografia: C 44%, M 50%, Y 68%, K 45% (~ Pantone 448C)
	\definecolor{cover-foreground}{cmyk}{0.44 0.50 0.68 0.45}
}
{
	%Cor do fundo da capa, lombada e contracapa: C 17%, M 26%, Y 67%, K 4% (~ Pantone 480C)
	\definecolor{cover-background}{cmyk}{0.17 0.26 0.67 0.04}
	%Cor das marcas e Tipografia: C 44%, M 50%, Y 68%, K 45% (~ Pantone 448C)
	\definecolor{cover-foreground}{cmyk}{0.44 0.50 0.68 0.45}
}



\usepackage[page,titletoc,title]{appendix} 

\usepackage{subcaption}

\usepackage[sort&compress,square]{natbib}
\renewcommand{\cite}[1]{\citep{#1}}

\usepackage[nottoc,notlot,notlof]{tocbibind}

\usepackage{notoccite}
\usepackage{titlesec}

%Remove all paragraph indentation
\setlength\parindent{0pt}


% --------------- headers and footers control ----------------------
\usepackage{fancyhdr}
\pagestyle{fancy}
\renewcommand{\chaptermark}[1]{\markboth{\MakeUppercase{#1}}{}}
\fancyhead[LE]{\ifnum \value{chapter}>0\sl{\nouppercase{\chaptertitlename~\thechapter}}\fi}
\fancyhead[RE]{}
\fancyhead[LO]{}
\fancyhead[RO]{\sl{\nouppercase{\leftmark}}}

%\fancyhead[RE,LO]{\rightmark}
\renewcommand{\headrulewidth}{0.4pt}
%---------------------------------------------------------------------


\makenoidxglossaries
\usepackage[hidelinks,bookmarks=true]{hyperref}


\hypersetup{
    pdftitle={\titleEN},
    pdfauthor={\author}
}

\thispagestyle{plain}

\newacronym{os}{OS}{Operating System}
\newacronym{nyi}{NYI}{Not Yet Implemented}
\newacronym{dei}{DEI}{Department of Informatics Engineering}


\usepackage{etoolbox}


\begin{document}
\setlength{\parskip}{0.7em}

%---------------------------------- FRONTMATTER ---------------------------------- 
\frontmatter

\newgeometry{left=2.5cm,right=2.5cm, bottom=4cm, top=5cm}

\thispagestyle{empty}
\pagecolor{cover-background}\afterpage{\nopagecolor}

\color{cover-foreground}

%\mbox{}
%vspace{0.5cm}

\begin{center}

\begin{figure}[h!]
	\begin{center}
	\includegraphics[height=5cm]{frontmatter/logo-uc-gold.pdf}	
	\end{center}
\end{figure}




\vspace{2cm}


\Large
\author

\vspace{1cm}

\LARGE

\ifthenelse{\equal{\docLanguage}{en}}
{\textsc{\textbf{\titleEN}}}
{\textsc{\textbf{\titlePT}}}



\Large
\vspace{-0.1cm}

\ifthenelse{\equal{\docLanguage}{en}}
{\textsc{\subtitleEN}}
{\textsc{\subtitlePT}}




%\vspace{3cm}


\begin{textblock*}{16cm}(2.5cm,19.5cm) % {block width} (coords) 
\color{cover-foreground}



\normalsize
\begin{spacing}{1.1}


%\textbf{

\ifthenelse{\equal{\docLanguage}{en}}
{
	\ifthenelse{\equal{\thesistype}{phd}}
	{\textbf{\descPhDEN}}
	{\textbf{\descMScEN}}
}
{
	\ifthenelse{\equal{\thesistype}{phd}}
	{\textbf{\descPhDPT}}
	{\textbf{\descMScPT}}
}


%}

\end{spacing}

\end{textblock*}




%\vspace{1.5cm}

\large

\vspace*{\fill}
\color{cover-foreground}
\ifthenelse{\equal{\docLanguage}{en}}
{\dateEN}
{\datePT}




\end{center}


\restoregeometry
\color{black}

\cleardoublepage

\ifthenelse{\equal{\docLanguage}{en}}
{ 
	\newgeometry{left=3.5cm,right=2.5cm, bottom=3.5cm, top=3.5cm}

\thispagestyle{empty}




%\mbox{}
%vspace{0.5cm}

\begin{center}

\begin{figure}[h!]
	\begin{center}
	\includegraphics[height=7.5cm]{frontmatter/logo-dei.pdf}	
	\end{center}
\end{figure}



\vspace{1cm}


\Large
\author

\vspace{1cm}

\LARGE

\ifthenelse{\equal{\docLanguage}{en}}
{\textsc{\textbf{\titleEN}}}
{\textsc{\textbf{\titlePT}}}



\Large
\vspace{-0.1cm}

\ifthenelse{\equal{\docLanguage}{en}}
{\textsc{\subtitleEN}}
{\textsc{\subtitlePT}}







\begin{textblock*}{15cm}(3.5cm,21cm) % {block width} (coords) 




\normalsize
\begin{spacing}{1.1}


%\textbf{

\ifthenelse{\equal{\docLanguage}{en}}
{
	\ifthenelse{\equal{\thesistype}{phd}}
	{\textbf{\titlePageDescPhdEN}}
	{\textbf{\titlePageDescMScEN}}
}
{
	\ifthenelse{\equal{\thesistype}{phd}}
	{\textbf{\titlePageDescPhdPT}}
	{\textbf{\titlePageDescMScPT}}
}


%}

\end{spacing}

\end{textblock*}




%\vspace{1.5cm}

\large

\vspace*{\fill}

\ifthenelse{\equal{\docLanguage}{en}}
{\dateEN}
{\datePT}




\end{center}


\restoregeometry


	\cleardoublepage	
	\def\docLanguage{pt}
	\newgeometry{left=3.5cm,right=2.5cm, bottom=3.5cm, top=3.5cm}

\thispagestyle{empty}




%\mbox{}
%vspace{0.5cm}

\begin{center}

\begin{figure}[h!]
	\begin{center}
	\includegraphics[height=7.5cm]{frontmatter/logo-dei.pdf}	
	\end{center}
\end{figure}



\vspace{1cm}


\Large
\author

\vspace{1cm}

\LARGE

\ifthenelse{\equal{\docLanguage}{en}}
{\textsc{\textbf{\titleEN}}}
{\textsc{\textbf{\titlePT}}}



\Large
\vspace{-0.1cm}

\ifthenelse{\equal{\docLanguage}{en}}
{\textsc{\subtitleEN}}
{\textsc{\subtitlePT}}







\begin{textblock*}{15cm}(3.5cm,21cm) % {block width} (coords) 




\normalsize
\begin{spacing}{1.1}


%\textbf{

\ifthenelse{\equal{\docLanguage}{en}}
{
	\ifthenelse{\equal{\thesistype}{phd}}
	{\textbf{\titlePageDescPhdEN}}
	{\textbf{\titlePageDescMScEN}}
}
{
	\ifthenelse{\equal{\thesistype}{phd}}
	{\textbf{\titlePageDescPhdPT}}
	{\textbf{\titlePageDescMScPT}}
}


%}

\end{spacing}

\end{textblock*}




%\vspace{1.5cm}

\large

\vspace*{\fill}

\ifthenelse{\equal{\docLanguage}{en}}
{\dateEN}
{\datePT}




\end{center}


\restoregeometry


	\cleardoublepage
	\def\docLanguage{en}
}
{
	\newgeometry{left=3.5cm,right=2.5cm, bottom=3.5cm, top=3.5cm}

\thispagestyle{empty}




%\mbox{}
%vspace{0.5cm}

\begin{center}

\begin{figure}[h!]
	\begin{center}
	\includegraphics[height=7.5cm]{frontmatter/logo-dei.pdf}	
	\end{center}
\end{figure}



\vspace{1cm}


\Large
\author

\vspace{1cm}

\LARGE

\ifthenelse{\equal{\docLanguage}{en}}
{\textsc{\textbf{\titleEN}}}
{\textsc{\textbf{\titlePT}}}



\Large
\vspace{-0.1cm}

\ifthenelse{\equal{\docLanguage}{en}}
{\textsc{\subtitleEN}}
{\textsc{\subtitlePT}}







\begin{textblock*}{15cm}(3.5cm,21cm) % {block width} (coords) 




\normalsize
\begin{spacing}{1.1}


%\textbf{

\ifthenelse{\equal{\docLanguage}{en}}
{
	\ifthenelse{\equal{\thesistype}{phd}}
	{\textbf{\titlePageDescPhdEN}}
	{\textbf{\titlePageDescMScEN}}
}
{
	\ifthenelse{\equal{\thesistype}{phd}}
	{\textbf{\titlePageDescPhdPT}}
	{\textbf{\titlePageDescMScPT}}
}


%}

\end{spacing}

\end{textblock*}




%\vspace{1.5cm}

\large

\vspace*{\fill}

\ifthenelse{\equal{\docLanguage}{en}}
{\dateEN}
{\datePT}




\end{center}


\restoregeometry


	\cleardoublepage	
	\def\docLanguage{en}
	\newgeometry{left=3.5cm,right=2.5cm, bottom=3.5cm, top=3.5cm}

\thispagestyle{empty}




%\mbox{}
%vspace{0.5cm}

\begin{center}

\begin{figure}[h!]
	\begin{center}
	\includegraphics[height=7.5cm]{frontmatter/logo-dei.pdf}	
	\end{center}
\end{figure}



\vspace{1cm}


\Large
\author

\vspace{1cm}

\LARGE

\ifthenelse{\equal{\docLanguage}{en}}
{\textsc{\textbf{\titleEN}}}
{\textsc{\textbf{\titlePT}}}



\Large
\vspace{-0.1cm}

\ifthenelse{\equal{\docLanguage}{en}}
{\textsc{\subtitleEN}}
{\textsc{\subtitlePT}}







\begin{textblock*}{15cm}(3.5cm,21cm) % {block width} (coords) 




\normalsize
\begin{spacing}{1.1}


%\textbf{

\ifthenelse{\equal{\docLanguage}{en}}
{
	\ifthenelse{\equal{\thesistype}{phd}}
	{\textbf{\titlePageDescPhdEN}}
	{\textbf{\titlePageDescMScEN}}
}
{
	\ifthenelse{\equal{\thesistype}{phd}}
	{\textbf{\titlePageDescPhdPT}}
	{\textbf{\titlePageDescMScPT}}
}


%}

\end{spacing}

\end{textblock*}




%\vspace{1.5cm}

\large

\vspace*{\fill}

\ifthenelse{\equal{\docLanguage}{en}}
{\dateEN}
{\datePT}




\end{center}


\restoregeometry


	\cleardoublepage
	\def\docLanguage{pt}
}


\ifthenelse{\equal{\thesistype}{phd}}
{ 
	\thispagestyle{plain}

\vspace*{\fill}

The work presented in this thesis was carried out within the Software and Systems Engineering (SSE) group of the Centre for Informatics and Systems of the University of Coimbra (CISUC) in the context of the following projects:

\begin{itemize}
	\item AIDA: Adaptive, Intelligent and Distributed Assurance Platform. FCT (CMU-PT) POCI-01-0247-FEDER-045907.
	\item AIDA: Adaptive, Intelligent and Distributed Assurance Platform. FCT (CMU-PT) POCI-01-0247-FEDER-045907.	
	\item Other...		
\end{itemize}

This work has been partially supported by the following grants...

\begin{itemize}
	\item Grant nr 1.
	\item Grant nr 2.	
	\item Grant nr ...	
\end{itemize}
\vspace*{\fill}

\begin{figure}[!hb]
		\centering
		\includegraphics[width=3cm]{frontmatter/sponsor1.pdf}
\end{figure}


	\cleardoublepage
}{}


\thispagestyle{plain}

\ifthenelse{\equal{\docLanguage}{en}}
{\section*{Acknowledgements}}
{\section*{Agradecimentos}}


I would like to thank...



\newpage
\thispagestyle{empty}

\begin{center}
\vspace*{\fill}

\ifthenelse{\equal{\docLanguage}{en}}{
  {\textsc{Statement on the use of artificial intelligence tools}\par}

  {\footnotesize
This document was reviewed and refined with the assistance of Large Language Models, such as ChatGPT, or similar tools, which helped check grammar, correct typos, and enhance clarity. Any Artificial Intelligence-generated text or content is clearly indicated within the document and is used only to a limited extent. The overall content and ideas remain solely the responsibility of the author.
}
}{
  {\textsc{Declaração sobre o uso de ferramentas de inteligência artificial}\par}

  {
\footnotesize
Este documento foi revisto e melhorado com o auxílio de Large Language Models, como o ChatGPT ou de ferramentas similares, que ajudaram a verificar a gramática, a corrigir gralhas e a aumentar a clareza. Qualquer texto ou conteúdo gerado por Inteligência Artificial é claramente indicado no documento e é utilizado apenas de forma limitada. O conteúdo e as ideias gerais são da exclusiva responsabilidade do autor.}
}

\vspace*{0pt}
\end{center}
\cleardoublepage




\ifthenelse{\equal{\docLanguage}{en}}
{
	\thispagestyle{plain}

\section*{Abstract}
\label{sec:abstract}


Lorem ipsum dolor sit amet, consectetur adipiscing elit. Vestibulum sed efficitur magna, et laoreet elit. Praesent sed nulla est. Cras in sem nec sem posuere hendrerit eget eu ligula. Duis tincidunt, lacus quis sollicitudin porttitor, nibh lorem tempor elit, nec tempus nunc ante id felis.




\section*{Keywords}
\label{sec:keywords}

Lorem ipsum dolor sit amet, consectetur adipiscing elit. Vestibulum sed efficitur magna, et laoreet elit. Praesent sed nulla est. Cras in sem nec sem posuere hendrerit eget eu ligula. Duis tincidunt, lacus quis sollicitudin porttitor, nibh lorem tempor elit, nec tempus nunc ante id felis. 

	\cleardoublepage
	\thispagestyle{plain}

\section*{Resumo}
\label{sec:resumo}


Lorem ipsum dolor sit amet, consectetur adipiscing elit. Vestibulum sed efficitur magna, et laoreet elit. Praesent sed nulla est. Cras in sem nec sem posuere hendrerit eget eu ligula. Duis tincidunt, lacus quis sollicitudin porttitor, nibh lorem tempor elit, nec tempus nunc ante id felis.




\section*{Palavras-Chave}
\label{sec:palavras}

Lorem ipsum dolor sit amet, consectetur adipiscing elit. Vestibulum sed efficitur magna, et laoreet elit. Praesent sed nulla est. Cras in sem nec sem posuere hendrerit eget eu ligula. Duis tincidunt, lacus quis sollicitudin porttitor, nibh lorem tempor elit, nec tempus nunc ante id felis. 
	\cleardoublepage
}
{%else
	\thispagestyle{plain}

\section*{Resumo}
\label{sec:resumo}


Lorem ipsum dolor sit amet, consectetur adipiscing elit. Vestibulum sed efficitur magna, et laoreet elit. Praesent sed nulla est. Cras in sem nec sem posuere hendrerit eget eu ligula. Duis tincidunt, lacus quis sollicitudin porttitor, nibh lorem tempor elit, nec tempus nunc ante id felis.




\section*{Palavras-Chave}
\label{sec:palavras}

Lorem ipsum dolor sit amet, consectetur adipiscing elit. Vestibulum sed efficitur magna, et laoreet elit. Praesent sed nulla est. Cras in sem nec sem posuere hendrerit eget eu ligula. Duis tincidunt, lacus quis sollicitudin porttitor, nibh lorem tempor elit, nec tempus nunc ante id felis. 
	\cleardoublepage
	\thispagestyle{plain}

\section*{Abstract}
\label{sec:abstract}


Lorem ipsum dolor sit amet, consectetur adipiscing elit. Vestibulum sed efficitur magna, et laoreet elit. Praesent sed nulla est. Cras in sem nec sem posuere hendrerit eget eu ligula. Duis tincidunt, lacus quis sollicitudin porttitor, nibh lorem tempor elit, nec tempus nunc ante id felis.




\section*{Keywords}
\label{sec:keywords}

Lorem ipsum dolor sit amet, consectetur adipiscing elit. Vestibulum sed efficitur magna, et laoreet elit. Praesent sed nulla est. Cras in sem nec sem posuere hendrerit eget eu ligula. Duis tincidunt, lacus quis sollicitudin porttitor, nibh lorem tempor elit, nec tempus nunc ante id felis. 

	\cleardoublepage
}

\ifthenelse{\equal{\thesistype}{phd}}
{
	\thispagestyle{plain}

The contributions of this thesis resulted in the following publications in international peer-reviewed conferences and journals:

\textbf{Journal papers:}
\begin{itemize}
	\item Publication1.
	\item Publication2.	
	\item Other...		
\end{itemize}

\textbf{Conference papers:}
\begin{itemize}
	\item Publication1.
	\item Publication2.	
	\item Other...		
\end{itemize}


\textbf{Workshop papers:}
\begin{itemize}
	\item Publication1.
	\item Publication2.	
	\item Other...		
\end{itemize}

The following papers are also related to this thesis, but were not included:

\textbf{Other papers:}
\begin{itemize}
	\item Publication1.
	\item Publication2.	
	\item Other...		
\end{itemize}
	\cleardoublepage
}{}
%-----------------------------------------------------------------------------------

% Set paragraph spacing to use in ToC and other lists
\setlength{\parskip}{0em}

% Make sure acronyms don't appear in the ToC
\glsunsetall
\tableofcontents
\glsresetall

\cleardoublepage


% List of acronyms
\glsnogroupskiptrue

\ifthenelse{\equal{\docLanguage}{en}}
{
	\def\acronymsTitle{Acronyms}
}
{
	\def\acronymsTitle{Acrónimos}
}
\printnoidxglossary[type=\acronymtype,style=list,title={\acronymsTitle}, nonumberlist]

\cleardoublepage

\listoffigures
\cleardoublepage

\listoftables
\cleardoublepage

%------------------------------
\mainmatter
\setlength{\parskip}{0.7em}
\chapter{Introduction}
\label{chap:introduction}

% PLAN: write in W11 (4--10 Jan), ~5--7 pages. Opening paragraph: the work is
% done in the scope of the MSc in Informatics Engineering (AI specialization)
% at the University of Coimbra; say what this chapter covers.

\section{Context and Motivation}
\label{sec:intro-context}

% - Automated decisions in high-stakes domains (credit, hiring, justice, health).
% - Explaining a negative decision is not enough: people need actionable
%   recourse \cite{wachter2018counterfactual, ustun2019actionable, karimi2022survey}.
% - Recourse itself can be unfair: cost and feasibility differ across groups
%   \cite{gupta2019equalizing, vonkugelgen2022fairness, barrainkua2026revisiting}.
% - Interpretable-by-design models as an alternative to post-hoc explanation
%   \cite{rudin2019stop}; KANs as a recent candidate \cite{liu2025kan}.
% - Close with the problem statement in one or two sentences.

\section{Objectives and Research Questions}
\label{sec:intro-objectives}

% - Main goal in one bold sentence (as in the example reports).
% - RQ1--RQ4: draft in W1 with the supervisor; final wording comes from
%   Section~\ref{sec:lr-gaps}.
% - Expected contributions.

\section{Work Plan and Risk Analysis}
\label{sec:intro-workplan}

\subsection{Work Plan}
\label{subsec:intro-workplan}

% - First semester: what was actually done vs. what was planned
%   (Gantt figure, planned vs. real bars, and a short note on deviations).
% - Second semester (February--June 2027): Gantt figure + monthly bullets.

\subsection{Risk Analysis and Mitigation}
\label{subsec:intro-risks}

% Table: risk statement | strategy (mitigate / accept) | mitigation action.
% Candidates: KAN training cost/instability; no ground-truth causal graph in
% real data; conflicting fairness notions; being scooped by a new preprint.

\section{Document Structure}
\label{sec:intro-structure}

% One paragraph describing Chapters~\ref{chap:background}--\ref{chap:conclusion}.

\chapter{Background Knowledge}
\label{chap:background}

% PLAN: written alongside the clusters (W2--W6), ~15--18 pages.
% Intro paragraph: this chapter introduces the concepts needed to follow
% Chapter~\ref{chap:literature-review}. Notation is fixed here and reused
% everywhere else.

\section{Supervised Learning Setting and Notation}
\label{sec:bg-notation}

% (W2) Input space, label, classifier, sensitive attribute, cost function
% between an instance and its counterfactual, set of feasible actions.
% Keep it short; this is the reference for all later formulas.

\section{Kolmogorov-Arnold Networks}
\label{sec:bg-kan}

\subsection{The Kolmogorov-Arnold Representation Theorem}
\label{subsec:bg-ka-theorem}

% (W3) Statement of the theorem \cite{kolmogorov1957representation, arnold1957functions};
% why it was long considered irrelevant for learning \cite{girosi1989representation}.

\subsection{Architecture and Training}
\label{subsec:bg-kan-architecture}

% (W3) Learnable univariate edge functions (B-splines + residual basis),
% grid extension, KAN layers vs. MLP layers \cite{liu2025kan};
% caveats from fair comparisons \cite{yu2024kan}. Figure: KAN vs. MLP schematic.

\subsection{Interpretability: Pruning and Symbolic Regression}
\label{subsec:bg-kan-interpretability}

% (W3) Sparsification, pruning, symbolic formula extraction
% \cite{liu2025kan, liu2024kan}; relation to additive models
% \cite{agarwal2021neural}.

\section{Algorithmic Fairness}
\label{sec:bg-fairness}

\subsection{Sources of Bias}
\label{subsec:bg-bias-sources}

% (W5) Taxonomy of bias sources \cite{mehrabi2021survey, barocas2023fairness}.

\subsection{Group Fairness Criteria}
\label{subsec:bg-group-fairness}

% (W5) Demographic parity, equal opportunity, equalized odds, calibration
% \cite{hardt2016equality}; impossibility results
% \cite{kleinberg2017inherent, chouldechova2017fair}.

\subsection{Individual and Causal Fairness}
\label{subsec:bg-individual-causal-fairness}

% (W5) Fairness through awareness \cite{dwork2012fairness}; counterfactual
% fairness \cite{kusner2017counterfactual}; causal reasoning about
% discrimination \cite{kilbertus2017avoiding}.

\subsection{Bias Mitigation Strategies}
\label{subsec:bg-mitigation}

% (W5) Pre-, in- and post-processing \cite{caton2024fairness};
% fair representations \cite{zemel2013learning}.

\section{Causal Inference Foundations}
\label{sec:bg-causality}

\subsection{Structural Causal Models}
\label{subsec:bg-scm}

% (W6) SCMs, causal graphs, structural equations, exogenous noise
% \cite{pearl2009causality, peters2017elements}.

\subsection{Interventions and Counterfactuals}
\label{subsec:bg-counterfactuals}

% (W6) Ladder of causation; do-operator; abduction--action--prediction.
% Figure: a small SCM with an intervention (reuse it in Section~\ref{sec:bg-recourse}).

\section{Counterfactual Explanations and Algorithmic Recourse}
\label{sec:bg-recourse}

\subsection{Counterfactual Explanations}
\label{subsec:bg-ce}

% (W4) Optimisation formulation \cite{wachter2018counterfactual}; desiderata.

\subsection{From Counterfactual Explanations to Recourse}
\label{subsec:bg-ce-to-recourse}

% (W4) Actionability and cost \cite{ustun2019actionable}; recourse as minimal
% interventions instead of nearest counterfactuals \cite{karimi2021algorithmic}.

\subsection{Evaluation Criteria}
\label{subsec:bg-recourse-metrics}

% (W4) Validity, proximity, sparsity, plausibility, diversity, cost, runtime
% \cite{verma2024counterfactual, guidotti2024counterfactual}.

\section{Summary}
\label{sec:bg-summary}

% Half a page: how the four building blocks connect to the thesis.

\chapter{Literature Review}
\label{chap:literature-review}

% PLAN: the core of the January delivery, ~25--32 pages (W3--W8).
% Intro paragraph: positioning at the intersection of recourse, fairness and
% KANs, with causal modelling as the enabler (positioning figure). Explain that
% the chapter ends with a decision about the role of the KAN (decision model,
% causal model or recourse generator), taken up in Chapter~\ref{chap:proposed-approach}.
% Rule: every section ends with ``Summary and Implications'' tied to the thesis.

\section{Review Methodology}
\label{sec:lr-methodology}

% (W3) Protocol following \cite{kitchenham2007guidelines}; reporting following
% \cite{page2021prisma}; snowballing following \cite{wohlin2014guidelines}.

\subsection{Review Questions}
\label{subsec:lr-review-questions}

% The review's own questions (one per cluster), not the thesis RQs.

\subsection{Search Strategy}
\label{subsec:lr-search}

% Databases, preprint servers, venues, strings S1--S5 (table), search dates.

\subsection{Inclusion and Exclusion Criteria}
\label{subsec:lr-criteria}

% IC1--IC4, EC1--EC5, time windows.

\subsection{Study Selection}
\label{subsec:lr-selection}

% PRISMA flow diagram with the logged counts (finish in W8).

\section{Counterfactual Explanation and Recourse Generation Methods}
\label{sec:lr-recourse-methods}

% (W4) Taxonomy figure + comparison table:
% method | year | model access | causal? | actionability | plausibility | data | code.
% Seeds: \cite{karimi2022survey, verma2024counterfactual, guidotti2024counterfactual}.

\subsection{Optimisation-Based Methods}
\label{subsec:lr-optimisation}

% \cite{wachter2018counterfactual, ustun2019actionable, mothilal2020explaining}.

\subsection{Data-Manifold and Generative Methods}
\label{subsec:lr-generative}

% \cite{poyiadzi2020face, joshi2019towards, pawelczyk2020learning}.
% Note amortized generators (relevant to option O3).

\subsection{Causality-Aware Recourse}
\label{subsec:lr-causal-recourse}

% \cite{mahajan2019preserving, karimi2021algorithmic}.

\subsection{Robustness and Global Recourse}
\label{subsec:lr-robust-global}

% \cite{upadhyay2021towards, dominguezolmedo2022adversarial, rawal2020beyond};
% critical perspectives \cite{venkatasubramanian2020philosophical, barocas2020hidden}.

\subsection{Summary and Implications}
\label{subsec:lr-recourse-summary}

\section{Fairness of Algorithmic Recourse}
\label{sec:lr-fair-recourse}

% (W5) Comparison table:
% notion | definition | needs SCM? | mitigation proposed? | datasets.

\subsection{Group-Level Notions of Recourse Fairness}
\label{subsec:lr-group-recourse-fairness}

% Cost parity and burden \cite{gupta2019equalizing, sharma2020certifai};
% social burden \cite{barrainkua2026revisiting}; substantive equality
% \cite{bell2024fairness}.

\subsection{Individual and Causal Notions of Recourse Fairness}
\label{subsec:lr-causal-recourse-fairness}

% \cite{vonkugelgen2022fairness, ehyaei2023robustness}.

\subsection{Auditing and Mitigating Unfair Recourse}
\label{subsec:lr-recourse-mitigation}

% Subgroup audits \cite{kavouras2023fairness}; in-training vs. post-hoc
% mitigation (relevant to option O1).

\subsection{Summary and Implications}
\label{subsec:lr-fair-recourse-summary}

\section{Causal Modelling for Counterfactual Inference}
\label{sec:lr-causal-models}

% (W6) Comparison table:
% model | family | identifiability | mixed data? | interpretable mechanisms?

\subsection{Deep Structural Causal Models}
\label{subsec:lr-dscm}

% \cite{pawlowski2020deep, xia2023neural}.

\subsection{Flow-, Autoencoder- and Diffusion-Based Causal Models}
\label{subsec:lr-causal-generative}

% \cite{khemakhem2021causal, sanchezmartin2022vaca, javaloy2023causal, chao2023interventional}.

\subsection{Recourse under Imperfect Causal Knowledge}
\label{subsec:lr-imperfect-causal}

% \cite{karimi2020algorithmic}. Relevant to option O2: where does the graph come from?

\subsection{Summary and Implications}
\label{subsec:lr-causal-summary}

\section{Kolmogorov-Arnold Networks in Trustworthy Machine Learning}
\label{sec:lr-kan}

% (W7) Comparison table:
% variant/work | basis function | claim | tabular evaluation? | fairness/XAI/causal use?
% Surveys: \cite{somvanshi2024survey, hou2024comprehensive}.

\subsection{Variants and Efficiency}
\label{subsec:lr-kan-variants}

% \cite{li2024kolmogorov, yu2024kan} + variants found by the search.

\subsection{KANs on Tabular Data}
\label{subsec:lr-kan-tabular}

% \cite{poeta2024benchmarking, yadav2025kolmogorov}.

\subsection{Interpretability Claims and Evidence}
\label{subsec:lr-kan-interpretability}

% \cite{liu2025kan, liu2024kan}; contrast with \cite{rudin2019stop, agarwal2021neural}.

\subsection{KANs for Fairness, Explainability and Causality}
\label{subsec:lr-kan-trustworthy}

% \cite{priyadarshini2026learning} + works found by S5. State the gap
% (no KAN-based recourse work as of the search freeze), backed by the search log.

\subsection{Summary and Implications}
\label{subsec:lr-kan-summary}

\section{Datasets, Metrics and Benchmarking Practice}
\label{sec:lr-datasets}

% (W8) Datasets table:
% dataset | domain | n | features | sensitive attribute | actionable features | causal graph? | used by.
% \cite{lequy2022survey, becker1996adult, hofmann1994statlog, angwin2016machine,
%       wightman1998lsac, yeh2009comparisons}; benchmarking \cite{pawelczyk2021carla}.
% Also: which metrics and statistical tests are reported.

\section{Synthesis and Research Gaps}
\label{sec:lr-gaps}

% (W8) 1--2 pages. Cross-cluster synthesis -> explicit gaps -> final RQs
% (copy them back to Section~\ref{sec:intro-objectives}). End with the decision
% matrix for the KAN's role (options O1--O4).

\chapter{Proposed Approach}
\label{chap:proposed-approach}

% PLAN: draft in W9--W10 (holiday buffer), ~5--8 pages. Provisional by design:
% the intermediate report must show a justified direction, not final results.

\section{Problem Formulation}
\label{sec:pa-formulation}

% Formal fair recourse problem using the notation of Section~\ref{sec:bg-notation}:
% classifier, cost, actions, fairness constraint/objective on recourse.

\section{Candidate Architectures}
\label{sec:pa-architectures}

% Decision matrix from Section~\ref{sec:lr-gaps}:
% O1 KAN as decision model | O2 KAN as causal model | O3 KAN as recourse
% generator | O4 combination. Criteria: feasibility, novelty, evaluability,
% interpretability. State the recommended option.

\section{Experimental Design}
\label{sec:pa-experimental-design}

% Datasets table (from the Evaluation tab: \cite{lequy2022survey, becker1996adult,
% hofmann1994statlog, angwin2016machine, wightman1998lsac, yeh2009comparisons};
% benchmark \cite{pawelczyk2021carla}; tool \cite{wexler2020whatif}), baselines (MLP counterparts + recourse
% methods from Sections~\ref{sec:lr-cf-search} and~\ref{sec:lr-plausibility}), metrics
% (predictive, recourse, recourse-fairness), statistical tests, protocol.
% Optional: preliminary sanity experiment, if agreed with the supervisor.

\section{Summary}
\label{sec:pa-summary}

\chapter{Conclusion}
\label{chap:conclusion}

% PLAN: write in W11, ~1--2 pages.
% - What the first semester achieved (review scope, number of studies, main
%   findings per cluster).
% - The gaps identified and the chosen direction (Chapter~\ref{chap:proposed-approach}).
% - Next steps for the second semester (link to Section~\ref{subsec:intro-workplan}).


\bibliographystyle{plainnat} 
%\bibliographystyle{apalike}

\ifthenelse{\equal{\docLanguage}{en}}
{
	\renewcommand{\bibname}{References}
}
{
	\renewcommand{\bibname}{Referências}
}


\cleardoublepage

% Custom headers style for the references chapter
%--------------
\phantomsection
\renewcommand{\chaptermark}[1]{} % disable automatic chaptermark
\fancyhead[LE,RO]{\sl\MakeUppercase{\bibname}} 
%--------------
\bibliography{references}

\cleardoublepage
%reset fancyheaders definitions
%--------------
\renewcommand{\chaptermark}[1]{\markboth{\MakeUppercase{#1}}{}}
\fancyhead[LE]{\ifnum \value{chapter}>0\sl{\nouppercase{\chaptertitlename~\thechapter}}\fi}
\fancyhead[RE]{}
\fancyhead[LO]{}
\fancyhead[RO]{\sl{\nouppercase{\leftmark}}}
%--------------

\begin{appendices}



	

\chapter{Title for the first appendix}
\label{appendixA}

Lorem ipsum dolor sit amet, consectetur adipiscing elit. Vestibulum sed efficitur magna, et laoreet elit. Praesent sed nulla est. Cras in sem nec sem posuere hendrerit eget eu ligula. Duis tincidunt, lacus quis sollicitudin porttitor, nibh lorem tempor elit, nec tempus nunc ante id felis.


\cleardoublepage
\chapter{Title for the second appendix}
\label{appendixB}

Lorem ipsum dolor sit amet, consectetur adipiscing elit. Vestibulum sed efficitur magna, et laoreet elit. Praesent sed nulla est. Cras in sem nec sem posuere hendrerit eget eu ligula. Duis tincidunt, lacus quis sollicitudin porttitor, nibh lorem tempor elit, nec tempus nunc ante id felis.


\end{appendices}
%------------------------------


\end{document}
%Document end
